%!TEX root = ../template.tex
%%%%%%%%%%%%%%%%%%%%%%%%%%%%%%%%%%%%%%%%%%%%%%%%%%%%%%%%%%%%%%%%%%%%
%% abstract-pt.tex
%% NOVA thesis document file
%%
%% Abstract in Portuguese
%%%%%%%%%%%%%%%%%%%%%%%%%%%%%%%%%%%%%%%%%%%%%%%%%%%%%%%%%%%%%%%%%%%%

\typeout{NT FILE abstract-pt.tex}%

A Teoria de Linguagens Formais e de Autómatos (FLAT) é rigorosa e frequentemente percecionada como abstrata, o que dificulta a aprendizagem inicial. 
Esta dissertação estende a biblioteca \textbf{OCamlFLAT} com suporte para \emph{gramáticas de atributos}, incluindo atributos sintetizados e herdados, 
estratégias de avaliação e verificações estáticas para detetar erros de especificação comuns. A contribuição centra-se num núcleo executável fiável e 
em interfaces claras.

O plano original incluía melhorias na aplicação web \textbf{OFLAT} e uma avaliação pedagógica. No âmbito final, o trabalho limita-se aos 
a estender a biblioteca \textbf{OCamlFLAT} com as estruturas necessárias para integração futura. As componentes gráficas e o estudo empírico do 
impacto pedagógico ficam \emph{fora do âmbito} e são delineadas como trabalho futuro.

A dissertação revê a fundamentação relevante sobre gramáticas de atributos no contexto da teoria da computação, realiza um levantamento das 
ferramentas educativas relacionadas e descreve a arquitetura e a implementação das extensões ao \textbf{OCamlFLAT}. Ao estabilizar o núcleo e 
definir pontos de integração, o trabalho fornece uma base para posterior \emph{visualização} e utilização em \emph{contexto letivo}.
% Palavras-chave do resumo em Português
% \begin{keywords}
% Palavra-chave 1, Palavra-chave 2, Palavra-chave 3, Palavra-chave 4
% \end{keywords}
\keywords{
  Teoria de Linguagens Formais \and
  Autómatos \and
  OCaml \and
  OFLAT \and
  OCamlFLAT \and
  Programação Funcional \and
  Gramáticas Formais
}
% to add an extra black line
